\documentclass[12pt]{article}
\usepackage[margin=1in]{geometry}
\usepackage{enumitem}
\usepackage{titlesec}
\usepackage{parskip}
\usepackage{amsmath}
\usepackage{amssymb}
\usepackage{amsfonts}
\usepackage{hyperref}
\usepackage{fontenc}
\usepackage{xcolor}

\title{\textbf{Synthetic Arguments}}

\begin{document}

\maketitle
\newpage



%======================
%General / Core
%======================






\section*{Collective Action}

\begin{itemize}
    \item Ziblatt 2017: can the elites solve a collective action problem in facilitating a transition to a coherent conservative political movement?
\end{itemize}




\section*{Development of Democracy}

\begin{itemize}
    \item Democracy as \emph{contestation}: there exists broad inclusion and a legitimate opposition with a legitimate chance of gaining power
    \begin{itemize}
        \item Dahl 1972: Democratization and development of the opposition are separate. Two dimensions: public contestation within politics, and the right to participate within politics. High contestation and high inclusiveness leads to polyarchy; other combinations are lesser
        \item Double-transition rule: a nation is considered a developed democracy when rule transitions $A \to B \to A$
        \item Dahl 1972: in the absence of any conflict, it is difficult to retain legitimate competition
    \end{itemize}
    \item Expansion of the franchise:
    \begin{itemize}
        \item Acemoglu \& Robinson 2000 -- the west extended the franchise as a credible commitment mechanism to ``buy off'' commoners
        \item North \& Weingast 1989 -- expansion of property rights in 18th century England as a means of commitment against exploitation and to secure public investment / credit
        \item Dahl 1972: three paths to polyarchy, (1) liberalization before inclusiveness, (2) inclusiveness before liberalization, (3) sudden shock (often unstable). Central question: do you make a hegemonic democracy first, or do you expand inclusion but retain authoritarianism under that broader franchise? Less common on the second one
    \end{itemize}
    \item Expansion of the franchise -- intentionally or unintentionally -- requires additional institutional development:
    \begin{itemize}
        \item Dahl 1972: Candidates and parties change as the franchise expands, in an effort to better achieve their electoral goals. Often results in new parties or party change (see Muller 1997 for a more comprehensive discussion of parties changing)
        \item Huntington 1968: without sufficient political development to keep up with economic development -- which itself will bring greater inclusivity -- we should expect discontent and instability
        \item Ziblatt 2017: If this development does not allow for conservative parties to be meaningfully included in government, then the transition to democracy is likely to be unstable and short-term
    \end{itemize}
    \item Levy 2006 -- transition of the goals of the democratic state away from market-steering to market-supporting
\end{itemize}




\section*{Historical Prerequisites for Democracy}

\begin{itemize}
    \item How agrarian classes developed:
    \begin{itemize}
        \item Dahl 1972: Farming societies differentiated into peasantry (hegemony) or free farmers (polyarchy)
    \end{itemize}
    \item Ziblatt 2017: whether or not the elites were able to organize into a coherent political party/structure that was meaningfully involved in the democratic process, such that their interests could still be represented to some extent
    \item Huntington 1968 ``Green uprising'': alliance of some sectors of rural and urban populations (see: workers/labour) that overcomes tensions between urbanites and rural elites and allows for stability
    \begin{itemize}
        \item Also Huntington: ``A society is vulnerable to revolution only when the opposition of the middle class to the political system coincides with the opposition of the peasants. Once the middle class becomes conservative, rural rebellion is still possible, b ut revolution is not.''
    \end{itemize}
\end{itemize}



\section*{Nation-building \& Nationality}


\begin{itemize}
    \item Four primary factors in whether an individual adopts new political beliefs: exposure, prestige, consistency w/ beliefs, consistency w/ experience \textcolor{red}{CITE}
    \item Cross-cutting cleavages:
    \begin{itemize}
        \item threaten polyarchy (Dahl 1972)
    \end{itemize}
\end{itemize}



\section*{Regime Transitions}

\begin{itemize}
    \item Ziblatt 2017: enduring stability of a democratic transition requires an organized conservative movement such that the elites are not existentially threatened
    \item Huntington 1968: regime transition as a byproduct of institutional development (that is, the type of regime changes)
    \item Huntington 1968 on Modernization: replacement of esoteric authority with legitimate justifications; rise of new goals and institutions to fulfill those goals; expansion of political participation
    \item Levy 2006 -- a little bit of a different lens, but change in the focus of democratic regimes
\end{itemize}





\section*{Backsliding/Populism}


\begin{itemize}
    \item Ziblatt 2017: the rise of populist movements is antithetical to the democratic need for an organized conservative movement and can subsequently lead to instability or backsliding
    \item Could use Huber \& Shipan 2002 to talk briefly about regulatory / bureaucratic capture, although it is more of an example of how to \emph{avoid} these outcomes via institutional design
\end{itemize}



\section*{International Influences}

\begin{itemize}
    \item Dahl 1972: Foreign influences can directly impose institutions, affect beliefs, or affect decision sets
    \item Gourevitch 1978: foreign politics have significant impacts on domestic structures BEYOND just the invasion effect
\end{itemize}




\section*{Party Politics}

\begin{itemize}
    \item Parties develop as coordination mechanisms to facilitate generally-common individual behaviors into tangible political action
    \item Cox 1997: The number of relevant parties or candidates we should expect is a function of both the number of relevant cleavages seeking representation and the permissiveness of electoral institutions to allow for such representation. Should expect $M+1$ relevant competitors where $M$ is district magnitude / number of seats available; Duverger is just this with $M=1$
    \begin{itemize}
        \item See Posner 2004 for why some cleavages are politically salient but others are not
        \item See Fearon \& Laitin 2003 for why some cleavages result in violence/instability but others do not
    \end{itemize}
\end{itemize}








%======================
% Theory of Science
%======================






\section*{Methodological Reasoning}






\section*{Empirical Logic \& Evaluation}

\begin{itemize}
    \item Huber \& Shipman 2002: use the institutional environment as a separate variable rather than a fixed component of situational context
\end{itemize}




\section*{Rationality \& Limits}

General thoughts: 

\begin{itemize}
    \item 
\end{itemize}

Specific examples:

\begin{itemize}
    \item Cox 1997: much of electoral politics happens outside of equilibrium
    \item Lijphart 1999: (my critique) characterizes a great deal of out-of-equilibrium behavior rather that ultimately proved short lived and ineffective
\end{itemize}



\section*{Positive Theory}



%======================
%Institutions
%======================



\section*{Foundations of Institutions}





\section*{Legislative Principal/Agents}

\begin{itemize}
    \item Cartel theory: parties are weak but legislators are strong; parties effectively cannot control their members (at least in US congress) and so instead must resort to influence and responsiveness
    \item Legislators face an inherent incentive to prioritize their own voters, but this can be counterproductive for their party -- subsequently, a foundational tension exists between legislators and their party leaders about whether to vote in line or pursue individual goals:
    \begin{itemize}
        \item Diermeier \& Vlaicu 2011: collective action problem where legislator can extract personal rents at the expense of party unity
        \item Jones \& Hwang 2005: Cartel theory seems to occur in Argentina
        \item Kats \& Mair 1995: Cartel theory; mass party is a single point in party development and cartels create a near-total overlap between the party and the state
        \item Crisp et al 2004: Legislators decide whether to focus on local or national issues
        \item Carey \& Shugart 1995: When voters decide to prioritize personal versus party reputation; how this might change across voting systems (big note here: closed list is hard-for-party)
        \item Cox 1987: How this power cessation to the party occurred (by means of solving an information problem)
    \end{itemize}
    \item Legislators have multiple principals to whom they are agents (Carey 2009) -- the more agents there are, and the more authority they have, and the more their goals conflict, the less effective we should expect legislative outcomes to be.
    \begin{itemize}
        \item Tsebelis 1999: the more veto players, the less effective we should expect the legislature to be -- laws passed will be less significant
    \end{itemize}
    \item For legislatures as principals, see Huber \& Shipan 2002
    \item Huber \& Shipan 2002: legislative capacity -- the costs that legislators face to drafting detailed legislation (opportunity, time, etc.)
\end{itemize}




\section*{Legislative Accountability \& Elections}

\begin{itemize}
    \item Cox 1997: Elections are market activities with a balanced supply and demand of candidates ($M+1$ rule) BUT can be broken by specifically ethnic or geographic parties, etc. (like SNP)
    \begin{itemize}
        \item Gibbard-Satterthwaite impossibility: either dictatorship, non-universality, not-strategy-proof, or indecisive
        \item Arrow impossibility: either dictatorship, or IIA, or non-pareto
    \end{itemize}
    \item Cox 1997: Three fundamental translations: social cleavages into partisan preferences, partisan preferences into votes, votes into seats
\end{itemize}



\section*{Majoritarian vs. Proportionalism}


\begin{itemize}
    \item Boix 1999: ruling parties can use the transition between majoritarianism and PR to maintain some exclusivity (even if possibly conceding contestation)
    \item Can tie to Ziblatt 2017 -- do the elites maintain meaningful political influence as democratization occurs? This is potentially more feasible under proportionalism than majoritarianism
    \item Huber \& Shipan 2002: Autonomy can also be largely influenced by whether or not the governing body has a majority versus a proportional lead: in the case of a majority, where ongoing conflict may be lessened (lesser need to involve the opposition), we can expect to see greater discretion to bureaucrats as they have fewer competing agents and are more trusted by the ruling government
\end{itemize}




\section*{Authoritarian Politics}

\begin{itemize}
    \item Repression is costly (Dahl 1972) and accordingly regimes without significant resources may just be more likely to tolerate an active opposition; conversely, if the state is highly endowed, this can raise the stakes of holding office and become a broader issue \textcolor{red}{cite this specifically? Svolik 2012?}
    \begin{itemize}
        \item Broad-scale oppression does not happen in polyarchies (Dahl 1972) -- low barriers to contestation protect inclusivity for marginalized groups and preemptively dissuades additional violence by the state
    \end{itemize}
\end{itemize}



\section*{Institutional Endogeneity}


\begin{itemize}
    \item Electoral institutions are created by political actors with incentives to maximize their chances of achieving, holding, and executing power (Boix 1999, Lijphart 1991 is a bit softer)
    \item The choice of electoral institutions tangibly impacts policy outcomes -- very similar to above point, with Cox's applied translations (social cleavages $\rightarrow$ partisan preferences $\rightarrow$ votes $\rightarrow$ seats) -- see Cox 1997 
    \item An abundance of resources makes holding power more valuable; Dahl 1972: those with the means to do so will attempt to bend political systems in their favor, but the distribution of those means itself is endogenous. (Also tied to inequality, which can arise in more developed economies)
    \item Huber \& Shipman 2002: bureaucratic institutions chosen based on the trust of the ruling party in bureacrats to implement their ideal policies
    \item Huntington 1968: while economic and social development may occur endogenously, political institutionalization requires applied efforts and organization
    \item Levy 2006: somewhat endogenous? Transition of institutional goals
\end{itemize}






\section*{Presidential vs Parliamentary}






\section*{Bureaucracy}


\begin{itemize}
    \item Huber \& Shipman 2002 -- all of deliberate discretion. Discretion to bureaucratic agencies is a function of the trust the ruling party has in its bureaucratic agents and the level of conflict/competition in place from other parties. Also note: ``Wise and salutary neglect'' as a quote from Edmund Burke 1775
    \item Informational asymmetry forms the basis of much of this discretion (Huber \& Shipman 2002) -- bureaucrats often have better information / more expertise than legislators, but also may have misaligned incentives
\end{itemize}



\section*{State Capacity}

\begin{itemize}
    \item General model of development in Huntington 1968
    \item Ziblatt 2017 -- capacity to involve the elites while still expanding the franchise?
    \item Ostrom 1990 -- is there enough social capital to make endogenous institutions effective? (Very much in line with Putnam here as well)
    \item Huber \& Shipan 2002 -- state capacity ultimately impacts how politicians will organize bureaucratic orders. Also, state capacity may be modeled as the costs that the state faces to governing (maybe this can be better viewed as efficiency?)
    \item Huntington 1968 (the whole thing, pretty much)
    \item Another point of state capacity: the political institutions must have the capacity to handle the involvement of all those who are part of the political sphere. If there is not sufficient institutional / state capacity to handle this, we should expect to see extrainstitutional (and often violent) alternatives
    \item Levy 2006: state capacity hasn't decreased, it's just shifted to alternative activities
\end{itemize}



\section*{Other assorted factors (and where to find them)}

\begin{itemize}
    \item Trust: Dahl 1972, mechanized through commitment devices by A\&R, N\&W, and within authoritarian states by Gandhi \& Przeworski 2007
\end{itemize}


\end{document}
